\chapter{Difference-in-Differences}\label{ch:b3:difference-in-differences}

On 1 April 1992 New Jersey raised its minimum wage from \usd{4.25} to
\usd{5.05} an hour. Its neighbour, Pennsylvania, did not. In February and
again in November of that year, two economists telephoned 410 fast-food
restaurants on both sides of the state line and asked how many people they
employed. % ledger: b3did:ck94.design
Employment per restaurant fell in Pennsylvania and rose slightly in New
Jersey. % ledger: b3did:ck94.t3
Neither change on its own says anything about the minimum wage: the season,
the regional economy and a hundred other things moved between February and
November. But those things moved on both sides of the line. The change in
Pennsylvania estimates what would have happened in New Jersey without the
rise, and the difference between the two changes estimates the effect of the
rise. This chapter is about that subtraction: when it works, how to look for
the cases where it does not, and what to do when the policy reaches different
places at different times, or only one place.

\begin{openingchart}{Mean full-time-equivalent employment per fast-food
restaurant, February and November 1992 (part-time workers count as one half).
Dashed: New Jersey's February level moved by Pennsylvania's change, the
path New Jersey would have followed under parallel trends.}{Card and Krueger
(1994), Table 3; the dashed line is our arithmetic on the printed means}
\begin{tikzpicture}
\begin{axis}[width=0.62\linewidth, height=0.42\linewidth, xmin=-0.3, xmax=1.7,
  xtick={0,1}, xticklabels={February 1992, November 1992}, ymin=17, ymax=25,
  ylabel={Employment per restaurant}, ymajorgrids, clip=false]
  \addplot[omDef, very thick, mark=*] table[x=x, y=pa] {figdata/bachelor-3/03-difference-in-differences/out/ck.csv}
    node[pos=1, anchor=south west, font=\footnotesize] {Pennsylvania};
  \addplot[omThm, very thick, mark=square*] table[x=x, y=nj] {figdata/bachelor-3/03-difference-in-differences/out/ck.csv}
    node[pos=1, anchor=north west, font=\footnotesize] {New Jersey};
  \addplot[omThm, thick, dashed, mark=o] table[x=x, y=cf] {figdata/bachelor-3/03-difference-in-differences/out/ck.csv}
    node[pos=1, right, font=\footnotesize, align=left] {New Jersey,\\ parallel path};
\end{axis}
\end{tikzpicture}
\end{openingchart}

Read the chart from the right. In November, New Jersey's restaurants employed
21.03 people on average. Had New Jersey followed Pennsylvania's path, which fell
by 2.16, it would have employed $20.44 - 2.16 = 18.28$. The difference, 2.75
workers per restaurant, is the estimate: more employment, not less, after the
rise. % figdata: ck_stats
(The paper prints 2.76, with a standard error of 1.36, computed from unrounded
means.) % ledger: b3did:ck94.t3
The whole estimate rests on the dashed line, which nobody observed. The
chapter's tools are ways of making that line believable, or of showing that
it is not. What the minimum wage does to employment is a question for the
chapter on the labour market; here the study serves only as the cleanest
picture of the design.

\begin{recall}
The Year 2 volume wrote causal questions with potential outcomes: $Y_i(1)$ and
$Y_i(0)$ are unit $i$'s outcome with and without treatment, only one of which is
observed, and the average treatment effect on the treated (ATT) is $\E[Y_i(1) -
Y_i(0) \mid D_i = 1]$. \Cref{ch:b3:panel-data} introduced unit and period
fixed effects and showed why standard errors must be clustered when errors and
treatment both persist within a unit.
\end{recall}

\begin{notation}[Symbols of this chapter]\label{not:b3:difference-in-differences:symbols}
$D_i = 1$ marks the treated group and $D_i = 0$ the comparison group; $t = 0$ is
before the treatment and $t = 1$ after. $Y_{it}(1)$ and $Y_{it}(0)$ are potential
outcomes in period $t$, and $\tau$ is the ATT in period 1 (not a tax rate). Event time, the number
of periods since a unit was first treated, is $k$.
\end{notation}

\section{Two groups, two periods}\label{sec:b3:difference-in-differences:two-by-two}

\begin{definition}[Difference-in-differences estimator]\label{def:b3:difference-in-differences:did}
\index{difference-in-differences estimator}
With a treated group observed before and after a treatment and a comparison
group observed at the same dates, the \emph{difference-in-differences
estimator} is
\[
  \hat\tau_{\mathrm{DiD}} = \bigl(\bar Y_{1,\mathrm{after}} - \bar Y_{1,\mathrm{before}}\bigr)
  - \bigl(\bar Y_{0,\mathrm{after}} - \bar Y_{0,\mathrm{before}}\bigr),
\]
where $\bar Y_{g,s}$ is the mean outcome of group $g$ in period $s$.
\end{definition}

\begin{definition}[Parallel-trends assumption]\label{def:b3:difference-in-differences:parallel-trends}
\index{parallel-trends assumption}
The \emph{parallel-trends assumption} is that, without the treatment, the
mean outcome of the treated group would have changed by the same amount as that
of the comparison group:
$\E[Y_{i1}(0) - Y_{i0}(0) \mid D_i = 1] = \E[Y_{i1}(0) - Y_{i0}(0) \mid D_i = 0]$.
\end{definition}

The assumption allows the two groups to differ in level --- Pennsylvania's
restaurants were larger than New Jersey's in February --- but not in the
change they would have had. It is weaker than random assignment, and it is
untestable for the period that matters, which lies after the treatment.

\begin{proposition}[Difference-in-differences identifies the ATT]\label{prop:b3:difference-in-differences:att}
Suppose the treated group is untreated in period 0 and treated in period 1, the
comparison group is untreated in both, and parallel trends hold. Then the
population difference in differences equals $\tau = \E[Y_{i1}(1) - Y_{i1}(0) \mid
D_i = 1]$.
\end{proposition}

\begin{proof}
Observed outcomes are $Y_{i1}(1)$ for the treated in period 1 and $Y_{it}(0)$ in
every other cell. The population difference in differences is therefore
\[
  \E[Y_{i1}(1) - Y_{i0}(0) \mid D_i = 1] - \E[Y_{i1}(0) - Y_{i0}(0) \mid D_i = 0].
\]
Add and subtract $\E[Y_{i1}(0) \mid D_i = 1]$ in the first term:
it becomes $\tau + \E[Y_{i1}(0) - Y_{i0}(0) \mid D_i = 1]$. By parallel trends the
second part equals the comparison group's change, which cancels.
\end{proof}

The proof shows what the estimator delivers when the assumption fails: $\tau$
plus the difference between the two groups' untreated trends. Every threat to
a difference-in-differences design is a reason why that difference might not
be zero.

\Cref{lst:b3:difference-in-differences:nsw} applies the estimator to a
job-training programme of the 1970s, the National Supported Work
demonstration, whose trainees are compared with a large sample of the
population survey. The programme ran in 1976--1977, so 1975 is before and 1978
after.

\omcode{code/bachelor-3/03-difference-in-differences/listings/did_nsw.py}{1}{27}{The difference in differences of mean earnings, then the same number as an OLS interaction.}\label{lst:b3:difference-in-differences:nsw}

\omoutput{code/bachelor-3/03-difference-in-differences/outputs/did_nsw.txt}{Mean earnings in dollars of the year.}\label{out:b3:difference-in-differences:nsw}

The trainees' earnings rose by \usd{4817}, the survey's by \usd{1196}, and the
estimate is \usd{3621}. % output: did_nsw
The survey sample earned nine times as much as the trainees in 1975; parallel
trends allow that, but it already suggests the two groups were on different
paths. The survey's earnings are also top-coded: incomes above a ceiling (about
\usd{25000} in these years) are recorded at the ceiling, which holds down the
survey's measured rise and so inflates the estimate.
% data: nsw('cps_controls') maxima 25,244 (1975) and 25,565 (1978), held by 3,007 and 4,145 people; psid_controls maxima above 120,000 That is part of why it is
twice the experimental answer; \cref{lst:b3:difference-in-differences:lalonde}
repeats the comparison with the second survey, whose earnings are not top-coded.

\section{Parallel trends, and how to look at them}\label{sec:b3:difference-in-differences:parallel}

Two kinds of evidence bear on the assumption. The first is the period before
the treatment: if the two groups moved in parallel for several years before,
it is more plausible that they would have done so after. The second is a
comparison where the answer is known to be zero.

\begin{definition}[Placebo test]\label{def:b3:difference-in-differences:placebo}
\index{placebo test}
A \emph{placebo test} applies an estimator to a period, a group or an outcome
that the treatment cannot have affected; an estimate far from zero reveals a
failure of the design's assumptions.
\end{definition}

The NSW data allow a direct test of the estimator, because the programme was a
randomised experiment: applicants were assigned by lottery to the programme or
to a control group, and the difference in 1978 earnings between the two
estimates the effect without any assumption about trends.

\begin{record}[LaLonde's test of non-experimental estimators]
LaLonde (1986) discarded the experiment's control group and replaced it with
samples from two large surveys, as an evaluator without an experiment would.
He then asked whether the estimators of the day --- regression with controls,
differences in differences, selection models --- recovered the experimental
answer. Most did not, and the estimates varied widely with the comparison
sample and the specification. Dehejia and Wahba (1999) showed that on a subset
of his sample, with two years of earlier earnings, careful matching came close
to the benchmark; Smith and Todd (2005) showed that this success depended on
the subset. The lesson has held: an estimator that works in one data set can
fail in the next, and the experimental benchmark is the only way to know.
\end{record}

\omcode{code/bachelor-3/03-difference-in-differences/listings/lalonde_benchmark.py}{1}{25}{The experimental benchmark against four estimators, with each survey as the comparison group.}\label{lst:b3:difference-in-differences:lalonde}

\omoutput{code/bachelor-3/03-difference-in-differences/outputs/lalonde_benchmark.txt}{Effects on 1978 earnings, dollars; robust standard errors in parentheses.}\label{out:b3:difference-in-differences:lalonde}

The experiment says \usd{1794}. % output: lalonde_benchmark
A raw comparison with either survey says the programme cost its trainees eight
to fifteen thousand dollars: the survey samples were simply far richer. The
difference in differences removes that level gap and lands on the right side
of zero, but at twice the benchmark with the population survey. Adding controls
moves it closer, to \usd{2264}; regression with two years of earlier earnings
lands below the benchmark, at \usd{794}. % output: lalonde_benchmark
No one could have chosen among these numbers without the experiment.

One reason is a pattern Ashenfelter (1978) found in training data: the
earnings of people who enter a programme fall in the year or two before they
enter, which is often why they enter. Measured from the bottom of that dip,
their earnings would have recovered anyway, and the recovery is counted as the
programme's effect. The trainees' earnings fell from 1974 to 1975, a little
more than the surveys' (\cref{exo:b3:difference-in-differences:6}). The test
sees only part of the dip, because the programme selected people who were out
of work when they applied.

\begin{model}{Difference-in-differences}
\assumes Without the treatment, the treated and comparison groups would have
changed by the same amount (parallel trends); the treatment does not affect
the comparison group; nothing else that changes the outcome happens to one
group only at the same date; the composition of each group is stable.
\predicts The treated group's outcome departs from the comparison group's
path at the treatment date, not before; placebo dates, groups and outcomes show
no effect.
\tested Against experiments: in the NSW data, the estimate with a survey
comparison group is twice the experimental one; with controls, closer. Its
pre-period predictions are checked in every careful application by an \omterm{def:b3:difference-in-differences:event-study}{event
study}.
\breaks When units are selected for treatment because of a temporary shock
(Ashenfelter's dip), when the treated group was already on a different trend,
when the treatment spills over into the comparison group, and, in the
two-way regression, when the treatment starts at different dates and its
effect changes over time (\cref{sec:b3:difference-in-differences:staggered}).
\end{model}

\section{The regression form and its errors}\label{sec:b3:difference-in-differences:regression}

\begin{proposition}[Difference-in-differences as a regression]\label{prop:b3:difference-in-differences:regression}
In the regression $Y_{it} = \kappa + \gamma D_i + \phi\, t + \tau D_i t + u_{it}$
over the four cells, with $t \in \{0, 1\}$, the OLS estimate of $\tau$ equals
the \omterm{def:b3:difference-in-differences:did}{difference-in-differences estimator}.
\end{proposition}

\begin{proof}
The regression has four parameters and the data four cells, so OLS fits each
cell mean exactly: $\bar Y_{0,0} = \kappa$, $\bar Y_{1,0} = \kappa + \gamma$,
$\bar Y_{0,1} = \kappa + \phi$ and $\bar Y_{1,1} = \kappa + \gamma +
\phi + \tau$. Then $(\bar Y_{1,1} - \bar Y_{1,0}) - (\bar Y_{0,1} - \bar
Y_{0,0}) = \tau$.
\end{proof}

The regression form earns its
place by what it adds: controls, more periods and groups, and standard errors.
The last need care. In the NSW listing each person appears twice, so the
standard error is clustered by person. In policy applications the treatment is
usually assigned to a few groups --- two states, a handful of provinces --- and
the errors of all units in a group share the group's shocks. Clustering by group
is then the rule (\cref{sec:b3:panel-data:clustering}), and with two groups it
is impossible: the comparison of one state with one other has, in the end, one
degree of freedom for the group shocks, however many restaurants were called.

\section{Event studies}\label{sec:b3:difference-in-differences:event-study}

\begin{definition}[Event study]\label{def:b3:difference-in-differences:event-study}
\index{event study}
An \emph{event study} regression, for units first treated at different dates,
replaces the single treatment dummy with one dummy for each event time $k$
(periods since the unit's first treatment), alongside unit and period effects,
and omits one reference period, usually $k = -1$. The coefficients on $k < 0$
(the leads) measure differences before the treatment; those on $k \ge 0$ (the
lags), the effect $k$ periods after.
\end{definition}

An \omterm{def:b3:difference-in-differences:event-study}{event study} plots the leads and lags with their confidence intervals. The
leads are a \omterm{def:b3:difference-in-differences:placebo}{placebo test} of parallel trends in the periods where it can be
checked, and the lags show how the effect builds up or fades. The chart shows
one on simulated data, in which half the treated units adopt in period 6 and
half in period 11 and a third group never does.

\begin{omchart}{Event study on a simulated staggered roll-out: estimated
coefficient for each period relative to adoption, with 95~\% confidence
intervals clustered by unit (reference period $k = -1$). Dashed: the true
effect, $0.5 + 0.1k$ from adoption on. 300 units, 16 periods.}{simulated data,
the chapter's seeded script}
\begin{tikzpicture}
\begin{axis}[omtime, width=0.8\linewidth, height=0.4\linewidth, xmin=-5.5, xmax=5.5,
  xtick={-5,...,5}, ymin=-0.3, ymax=1.2, xlabel={Periods since adoption, $k$},
  ylabel={Effect on the outcome}, legend pos=north west,
  y tick label style={/pgf/number format/.cd, fixed, precision=1}]
  \addplot[guide, domain=-5.5:5.5, samples=2, forget plot] {0};
  \draw[guide] (axis cs:-0.5,-0.3) -- (axis cs:-0.5,1.2);
  \addplot[omDef, only marks, mark=*, mark size=1.6pt,
    error bars/.cd, y dir=both, y explicit]
    table[x=k, y=est, y error plus expr=\thisrow{hi}-\thisrow{est},
      y error minus expr=\thisrow{est}-\thisrow{lo}]
    {figdata/bachelor-3/03-difference-in-differences/out/event_study.csv};
  \addplot[omThm, only marks, mark=diamond*, mark size=2pt] table[x=k, y=true]
    {figdata/bachelor-3/03-difference-in-differences/out/event_study.csv};
  \legend{estimate, true effect}
\end{axis}
\end{tikzpicture}
\end{omchart}

The four leads lie within a few hundredths of zero, and every interval covers
zero; the lags climb with the true effect. % figdata: event_study
Flat leads do not prove parallel trends after the treatment, and a test with
little power will find flat leads in data where the trends differ; but rising
or falling leads are hard evidence against the design. Two cautions. If units
react before the treatment --- shops stocking up before a tax rise is announced
to take effect --- the reference period $k = -1$ is itself treated and every
coefficient is shifted (\cref{exo:b3:difference-in-differences:10}). And the
regression that produced this chart is safe only because the effect was the
same for both cohorts; the next section shows what happens when it is not.

\section{Staggered adoption}\label{sec:b3:difference-in-differences:staggered}

\begin{definition}[Staggered adoption]\label{def:b3:difference-in-differences:staggered}
\index{staggered adoption}
A treatment has \emph{staggered adoption} when units start it at different
dates and, once treated, stay treated. Units that start at the same date form
a cohort.
\end{definition}

For years the standard estimator for staggered designs was the \omterm{def:b3:panel-data:twfe}{two-way fixed
effects} regression of \cref{ch:b3:panel-data}, $Y_{it} = a_i + \mu_t + \beta
D_{it} + u_{it}$, with $D_{it} = 1$ once unit $i$ is treated. Its coefficient
looks like a difference in differences, and with one adoption date it is one.
With several, it is a mixture.

\begin{proposition}[Two-way fixed effects as a weighted average]\label{prop:b3:difference-in-differences:bacon}
In a balanced panel with \omterm{def:b3:difference-in-differences:staggered}{staggered adoption}, the \omterm{def:b3:panel-data:twfe}{two-way fixed effects}
coefficient $\hat\beta$ is a weighted average, with non-negative weights summing
to one, of all the two-group, two-period difference-in-differences estimates
that can be formed from the data: each cohort against the never-treated; each
earlier cohort against a later one before the later one is treated; and each
later cohort against an earlier one after the earlier one is treated. The
weights grow with the size of the groups and with the variance of the
treatment dummy in each comparison (Goodman-Bacon 2021).
\end{proposition}

\admitted
The proof rewrites the \omterm{def:b3:panel-data:twfe}{two-way fixed effects} coefficient, by the
Frisch--Waugh--Lovell theorem, as a ratio of sums over pairs of cohorts; it
is in Goodman-Bacon's paper.

The last kind of comparison is the trouble. When the later cohort is treated,
the earlier one serves as its comparison, but the earlier one is itself still
being treated, and its effect may still be growing. Its change over that
period contains its own treatment effect, which is subtracted.

\begin{omfigure}
\begin{tikzpicture}[econ, xscale=0.75, yscale=0.6]
  \foreach \t in {1,...,12} \node[font=\scriptsize] at (\t,-0.6) {\t};
  \node[font=\footnotesize, anchor=east] at (0.4,2) {early cohort};
  \node[font=\footnotesize, anchor=east] at (0.4,1) {late cohort};
  \node[font=\footnotesize, anchor=east] at (0.4,0) {never treated};
  \fill[omDef!15] (0.5,1.7) rectangle (3.5,2.3);
  \fill[omDef] (3.5,1.7) rectangle (12.5,2.3);
  \fill[omThm!15] (0.5,0.7) rectangle (7.5,1.3);
  \fill[omThm] (7.5,0.7) rectangle (12.5,1.3);
  \fill[black!12] (0.5,-0.3) rectangle (12.5,0.3);
  \draw[very thick, dashed] (7.5,0.5) rectangle (12.5,2.5);
  \node[font=\footnotesize, align=center, anchor=south] at (10,2.55)
    {forbidden comparison: the late cohort's change\\ minus the early cohort's,
     whose effect is still growing};
  \node[font=\scriptsize, white] at (8,2) {treated};
  \node[font=\scriptsize, white] at (10,1) {treated};
\end{tikzpicture}
\par\smallskip{\small \omterm{def:b3:difference-in-differences:staggered}{Staggered adoption} in \cref{lst:b3:difference-in-differences:bacon}:
the early cohort is treated from period 4, the late one from period 8 (dark
bars). In periods 8--12 \omterm{def:b3:panel-data:twfe}{two-way fixed effects} uses the early cohort, already
treated, as a comparison for the late one.}
\end{omfigure}

\omcode{code/bachelor-3/03-difference-in-differences/listings/bacon.py}{1}{33}{\omterm{def:b3:panel-data:twfe}{Two-way fixed effects} against clean comparisons, when the effect grows with exposure.}\label{lst:b3:difference-in-differences:bacon}

\omoutput{code/bachelor-3/03-difference-in-differences/outputs/bacon.txt}{True effect: 1 per period of exposure.}\label{out:b3:difference-in-differences:bacon}

The true average effect on the treated is 4.29; \omterm{def:b3:panel-data:twfe}{two-way fixed effects} reports
2.55, % output: bacon
an error of about 40~\% with no sampling noise to blame. Here every unit's effect is
positive and the estimate is merely too small; in other designs the bad
comparisons can drive it below zero (\cref{exo:b3:difference-in-differences:8}).
The repair is to make only clean comparisons: each cohort against units not yet
treated, from the cohort's last untreated period. The last line of the listing
uses the simplest clean comparison, each cohort against the never-treated units,
4.24. % output: bacon
Several estimators formalise this --- Callaway and Sant'Anna (2021) average the
cohort-by-period comparisons; de Chaisemartin and D'Haultf\oe uille (2020) use
the units whose treatment changes between consecutive periods --- and differ
mainly in which comparisons they allow and how they weight them. An \omterm{def:b3:difference-in-differences:event-study}{event study}
with a never-treated group and separate coefficients for each cohort is another
clean route.

\section{Synthetic control}\label{sec:b3:difference-in-differences:synthetic}

Sometimes only one unit is treated: one country reunifies, one state passes a
tobacco law. No group of untreated units need follow the treated one in
parallel, and choosing a single comparison by hand invites the choice that
gives the answer one wanted.

\begin{definition}[Synthetic control; donor pool]\label{def:b3:difference-in-differences:synthetic-control}
\index{synthetic control}\index{donor pool}
For one treated unit and a set of untreated units, the \emph{donor pool}, a
\emph{synthetic control} is a weighted average of the donors, with
non-negative weights summing to one, chosen so that its outcome before the
treatment tracks the treated unit's as closely as possible. The estimated
effect in each later period is the gap between the treated unit and its
synthetic control.
\end{definition}

\begin{proposition}[Synthetic weights]\label{prop:b3:difference-in-differences:synth}
Given the treated unit's pre-treatment outcomes $y_1 \in \mathbb R^{T_0}$ and
the donors' $Y_0 \in \mathbb R^{T_0 \times J}$, the synthetic weights solve
$\min_w \lVert y_1 - Y_0 w \rVert^2$ subject to $w_j \ge 0$ and $\sum_j w_j = 1$:
a least-squares problem over the simplex, whose solution exists and is a point
of the convex hull of the donors' paths closest to the treated path.
\end{proposition}

\admitted
Existence follows because the simplex is compact and the objective continuous;
the solution is unique when the donors' paths are affinely independent. In
practice the paths are often replaced by pre-period averages of the outcome and
of predictors, weighted by a second, outer optimisation; the inner problem has
the same form.

The constraint is the method's discipline: no extrapolation beyond the donors,
and weights that a reader can inspect. A regression with unrestricted weights
would fit the pre-period at least as well, with negative weights on some
countries and weights above one on others.

\begin{record}[Reunification and West German income]
Abadie, Diamond and Hainmueller (2015) built a synthetic West Germany from
sixteen other rich countries, matched on income per person and its predictors
before 1990. Five received weight: Austria 0.42, the United States 0.22,
Japan 0.16, Switzerland 0.11 and the Netherlands 0.09. % ledger: b3did:adh15.weights
After 1990 the two paths separated. Over 1990--2003, West Germany's GDP per
person was about \usd{1600} a year below its \omterm{def:b3:difference-in-differences:synthetic-control}{synthetic control} on average,
about 8~\% of its 1990 level, and by 2003 the synthetic West Germany was about
12~\% richer than the real one. % ledger: b3did:adh15.effect
(A 2025 erratum by the authors states the dollar figures in current
purchasing-power dollars rather than constant ones; the percentages are
unchanged.) Placebo runs that pretend each donor reunified in 1990 give
ratios of post- to pre-1990 gap all below West Germany's, which is about 16;
ranking first of 17 by luck has a chance of $1/17 = 0.059$. % ledger: b3did:adh15.placebo
\end{record}

\Cref{lst:b3:difference-in-differences:synth} runs the method on a simulated
region and twenty donors, all driven by two common factors with different
loadings, so that no single donor matches the treated region.

\omcode{code/bachelor-3/03-difference-in-differences/listings/synth_demo.py}{1}{30}{Weights, pre-period fit, effect, and the placebo rank.}\label{lst:b3:difference-in-differences:synth}

\omoutput{code/bachelor-3/03-difference-in-differences/outputs/synth_demo.txt}{True effect $-2$ from period 21.}\label{out:b3:difference-in-differences:synth}

\begin{omchart}{A simulated region and its synthetic control (four donors carry
most of the weight). The policy starts in period 21; the true effect is $-2$.
Twenty donors driven by two common factors.}{simulated data, the chapter's
seeded script}
\begin{tikzpicture}
\begin{axis}[omtime, width=0.8\linewidth, height=0.4\linewidth, xmin=1, xmax=30,
  xlabel={Period}, ylabel={Outcome}, ymin=7, ymax=16, legend pos=north west]
  \draw[guide] (axis cs:20.5,7) -- (axis cs:20.5,16);
  \addplot[omDef, very thick, no marks] table[x=t, y=treated] {figdata/bachelor-3/03-difference-in-differences/out/synth.csv};
  \addplot[omThm, thick, dashed, no marks] table[x=t, y=synthetic] {figdata/bachelor-3/03-difference-in-differences/out/synth.csv};
  \legend{treated region, synthetic control}
\end{axis}
\end{tikzpicture}
\end{omchart}

Before period 21 the two paths are hard to tell apart: the root mean squared
gap is 0.26. After, the mean gap is $-1.68$ against a true $-2$. % output: synth_demo
Inference uses placebos. Each donor in turn is treated as if it had been
treated, a \omterm{def:b3:difference-in-differences:synthetic-control}{synthetic control} is built for it from the others, and the ratio of
its post-period to its pre-period gap is computed. The treated region's ratio,
6.8, is the largest of the 21; % output: synth_demo
if the policy had done nothing, the chance that the true unit would rank first
by luck is $1/21$, about 0.05. With fewer donors, no result can be significant
at that level, whatever the gap.

\section{Exercises}\label{sec:b3:difference-in-differences:exercises}

\begin{exercise}[$\star$]\label{exo:b3:difference-in-differences:1}
Read the opening chart.
\begin{enumerate}
\item Compute each state's change from February to November and the
difference in differences from the printed means.
\item Compute New Jersey's change in percent of its February level.
\item State in one sentence what the dashed line assumes.
\end{enumerate}
\end{exercise}

\begin{exercise}[$\star$]\label{exo:b3:difference-in-differences:2}
Average hourly wages (in dalens) in two regions of Brenland were 20 and 24 in
2023, and 21 and 26.5 in 2025. Only the second region cut its payroll tax in
2024.
\begin{enumerate}
\item Compute the difference-in-differences estimate of the tax cut's effect on
wages.
\item What assumption does it need, and what would you look at to check it?
\end{enumerate}
\end{exercise}

\begin{exercise}[$\star$]\label{exo:b3:difference-in-differences:3}
In the event-study chart of \cref{sec:b3:difference-in-differences:event-study},
read the coefficient at $k = 2$ and its interval. What would you conclude if the
leads had been $-0.4$, $-0.3$, $-0.2$ and $-0.1$ at $k = -5, \dots, -2$?
\end{exercise}

\begin{exercise}[$\star$]\label{exo:b3:difference-in-differences:4}
Describe two patterns of pre-treatment data that contradict parallel trends,
and one that does not contradict it even though the two groups differ.
\end{exercise}

\begin{exercise}[$\star\star$]\label{exo:b3:difference-in-differences:5}
Mean sales in two groups of shops go from 100 to 120 (group A, treated) and from
50 to 60 (group B, untreated).
\begin{enumerate}
\item Compute the difference in differences in levels and in logs.
\item Show that parallel trends cannot hold both in levels and in logs here, and
explain how one would choose between them.
\end{enumerate}
\end{exercise}

\begin{exercise}[$\star\star$]\label{exo:b3:difference-in-differences:6}
Use the NSW data to compute a placebo difference in differences from 1974 to
1975, both before the programme, with each survey as the comparison group.
\begin{enumerate}
\item Report the trainees' change, each survey's change, and the two placebo
estimates with their standard errors.
\item What does the result say about parallel trends, and what can it not say?
\end{enumerate}
\end{exercise}

\begin{exercise}[$\star\star$]\label{exo:b3:difference-in-differences:7}
Workers enter a training programme after a year in which their earnings fell
for temporary reasons. The evaluator measures the change in their earnings from
that year to two years after the programme, against a comparison group without
the dip. In which direction is the difference in differences biased? How would
choosing an earlier ``before'' year help?
\end{exercise}

\begin{exercise}[$\star\star$]\label{exo:b3:difference-in-differences:8}
Two equal-sized groups are observed in periods 1, 2 and 3. Group E is treated
from period 2, group L from period 3; there is no never-treated group. E's
effect is $\tau$ in period 2 and $\tau + \gamma$ in period 3 (here $\gamma$ is
the growth of the effect); L's effect is $\tau$ in period 3.
\begin{enumerate}
\item Compute the double-demeaned treatment dummy in each of the six cells.
\item Show that the \omterm{def:b3:panel-data:twfe}{two-way fixed effects} coefficient equals $\tau - \gamma/2$.
\item Give values of $\tau$ and $\gamma$, both positive, for which the
coefficient is negative. Which cell receives a negative weight?
\end{enumerate}
\end{exercise}

\begin{exercise}[$\star\star$]\label{exo:b3:difference-in-differences:9}
Kestria raises the sales tax on fuel in its two western provinces only. Propose
a placebo group, a placebo outcome and a placebo date for a
difference-in-differences study of the effect on fuel sales, and say what
each would rule out.
\end{exercise}

\begin{exercise}[$\star\star\star$]\label{exo:b3:difference-in-differences:10}
Simulate the design of the event-study chart, but let the units react one period
before adoption, with an effect of 0.3 at $k = -1$.
\begin{enumerate}
\item Estimate the \omterm{def:b3:difference-in-differences:event-study}{event study} with reference period $k = -1$ and report the
coefficients at $k = -3, -2, 0, 1, 2$.
\item Re-estimate with reference period $k = -2$. Explain the difference.
\end{enumerate}
\end{exercise}

\begin{exercise}[$\star\star\star$]\label{exo:b3:difference-in-differences:11}
A treated region has outcomes 4 and 6 in two pre-treatment periods. Three donors
have $(2, 4)$, $(6, 6)$ and $(4, 10)$.
\begin{enumerate}
\item Find the synthetic-control weights, and check that they reproduce the
pre-treatment path exactly.
\item After the treatment, the region's outcome is 5 and the donors' are 3, 7
and 9. Estimate the effect.
\item Why does the method forbid negative weights?
\end{enumerate}
\end{exercise}

\begin{exercise}[$\star\star\star$]\label{exo:b3:difference-in-differences:12}
A state introduces health insurance for people aged 65 and over. A
difference in differences of the old in that state against the old in other
states may be biased by state-wide shocks that hit the old differently.
\begin{enumerate}
\item Write the triple-difference estimator that also uses people under 65 in
both kinds of state.
\item State the assumption it needs, and show that it is weaker than
parallel trends for the old alone.
\end{enumerate}
\end{exercise}

\section{Problem: The Kestrian Border Tax}\label{sec:b3:difference-in-differences:problem}

\begin{problem}[{Weekend problem --- the Kestrian border tax}]\label{pb:b3:difference-in-differences:border}
Kestria (invented) has twelve provinces. In 2019 it raised its sales tax in the
three western provinces W1--W3, which border Brenland, where shopping became
relatively cheaper; in 2021 it raised it in the three eastern provinces E1--E3;
the six northern provinces N1--N6 kept the old rate. A statistician has the log
annual sales of 50 shops in each province, 2014--2024, and computes the tables
below. Standard errors are clustered by province. % figdata: kestria_*

\smallskip
\noindent\begin{minipage}[t]{0.4\linewidth}
\centering\small Mean log sales\par\smallskip
\pgfplotstabletypeset[col sep=comma, string type,
  columns/group/.style={column name={Group}, column type=l},
  columns/y2018/.style={column name={2018}, string type=false, numeric type, fixed, precision=4, zerofill},
  columns/y2020/.style={column name={2020}, string type=false, numeric type, fixed, precision=4, zerofill},
  every head row/.style={after row=\hline}, every last row/.style={after row=\hline}]
  {figdata/bachelor-3/03-difference-in-differences/out/kestria_means.csv}

\medskip
Synthetic W1\par\smallskip
\pgfplotstabletypeset[col sep=comma,
  columns/donor/.style={column name={Donor}, string type, column type=l},
  columns/weight/.style={column name={Weight}, fixed, precision=4, zerofill},
  every head row/.style={after row=\hline}, every last row/.style={after row=\hline}]
  {figdata/bachelor-3/03-difference-in-differences/out/kestria_synth.csv}
\end{minipage}\hfill
\begin{minipage}[t]{0.27\linewidth}
\centering\small West against north\par\smallskip
\pgfplotstabletypeset[col sep=comma,
  columns/k/.style={column name={$k$}, int detect},
  columns/est/.style={column name={Coef.}, fixed, precision=4, zerofill},
  columns/se/.style={column name={s.e.}, fixed, precision=4, zerofill},
  every head row/.style={after row=\hline}, every last row/.style={after row=\hline}]
  {figdata/bachelor-3/03-difference-in-differences/out/kestria_event.csv}
\end{minipage}\hfill
\begin{minipage}[t]{0.27\linewidth}
\centering\small Gap, W1 minus synthetic\par\smallskip
\pgfplotstabletypeset[col sep=comma,
  columns/year/.style={column name={Year}, int detect, /pgf/number format/set thousands separator={}},
  columns/gap/.style={column name={Gap}, fixed, precision=4, zerofill},
  every head row/.style={after row=\hline}, every last row/.style={after row=\hline}]
  {figdata/bachelor-3/03-difference-in-differences/out/kestria_synth_gap.csv}
\end{minipage}

\smallskip
\begin{center}\small
\pgfplotstabletypeset[col sep=comma,
  columns/estimator/.style={column name={Estimator of the tax's effect}, string type, column type=l},
  columns/est/.style={column name={Estimate}, fixed, precision=4, zerofill},
  columns/se/.style={column name={s.e.}, fixed, precision=4, zerofill, empty cells with={---}},
  every head row/.style={after row=\hline}, every last row/.style={after row=\hline}]
  {figdata/bachelor-3/03-difference-in-differences/out/kestria_estimates.csv}
\end{center}
The \omterm{def:b3:difference-in-differences:event-study}{event study} uses the western and northern shops only, with shop and year
effects and $k = -1$ (2018) as reference. The ``clean'' estimator averages, over
every treated province and year, the province's change since its last untreated
year minus the north's change. The data are invented, so the true average effect
over treated shop-years is known: it is the last line.

\medskip\noindent\textbf{Part I --- two groups, two years.}
\begin{enumerate}
\item From the table of means, compute each group's change from 2018 to 2020
and the difference in differences.
\item Convert it into a percentage change in sales.
\item Give a 95~\% confidence interval with the clustered standard error.
\item How many clusters does the standard error rest on? Why is that a concern?
\item The west's mean was below the north's already in 2018. Does that threaten
the estimate?
\item Name one event in 2019--2020, other than the tax, that would bias this
estimate, and say in which direction.
\end{enumerate}

\medskip\noindent\textbf{Part II --- the event study.}
\begin{enumerate}[resume]
\item Compute the $t$-statistic of each lead. What do the leads suggest?
\item The coefficient at $k = 1$ equals the difference in differences of Part I.
Explain why.
\item Describe the path of the lags from $k = 0$ to $k = 5$, and give one
economic reason why an effect on cross-border shopping might grow over two or
three years.
\item The lag at $k = 4$ has a larger standard error than the others. Why might
estimates far from the event be noisier?
\item Suppose shoppers had learned of the tax in mid-2018 and stocked up before
it. What would happen to the coefficients, given the reference year?
\item Why does this \omterm{def:b3:difference-in-differences:event-study}{event study} leave out the eastern provinces?
\end{enumerate}

\medskip\noindent\textbf{Part III --- the roll-out to the east.}
\begin{enumerate}[resume]
\item Which comparisons does the \omterm{def:b3:panel-data:twfe}{two-way fixed effects} regression over all twelve
provinces make? Name the three kinds.
\item Which of them is a forbidden comparison, and why does it bias the estimate
toward zero here?
\item Compute the bias of \omterm{def:b3:panel-data:twfe}{two-way fixed effects} as a percentage of the true effect.
\item The clean estimator comes within 0.001 of the truth. Which comparisons
does it avoid?
\item Compare the standard errors of the two-by-two and two-way estimates.
Is the smaller one a reason to prefer \omterm{def:b3:panel-data:twfe}{two-way fixed effects}?
\item If the effect had been constant from the first year on, would \omterm{def:b3:panel-data:twfe}{two-way
fixed effects} have been biased? Explain.
\end{enumerate}

\medskip\noindent\textbf{Part IV --- one province.}
\begin{enumerate}[resume]
\item Check that the synthetic weights for W1 sum to one, and say how many
donors receive weight.
\item Compute the root mean squared gap over 2014--2018.
\item Compute the mean gap over 2021--2024 and compare it with the west's
event-study lags for the same years.
\item Give one reason why W1's effect may be larger than the average western
province's.
\item Compute the root mean squared gap over 2019--2024 and the ratio of the
post- to the pre-period value.
\item With six donors, what is the smallest $p$-value a \omterm{def:b3:difference-in-differences:placebo}{placebo test} of ranks
can give? Is a rank of 1 out of 7 strong evidence?
\item Why should the \omterm{def:b3:difference-in-differences:synthetic-control}{donor pool} exclude the eastern provinces?
\item State the results of the problem: the effect of the tax on the sales of
western shops by each of the four estimators, and which you would report.
\end{enumerate}
\end{problem}
